% design diagram with major components

% APIs and differences (extra features) from HuggingFace.  


To address these challenges, we introduce the HPC-FAIR, a framework packaging a set of machine learning components with user-friendly API for HPC users. %Its integration is as straightforward as importing any standard Python package with the command \mintinline{python}{import hpc-fair}.  
%Built on the foundation of \face APIs, HPC-FAIR leverages the existing strengths of the \face platform while adding extra functionalities specifically designed for HPC tasks. 
Our design process begins by examining typical ML workflows for HPC tasks and then presenting HPC-FAIR components and API details to support these workflows.

% \subsection{Basic Operations}
% \textbf{Model load}

% \face provides convenient access to a vast array of pre-trained models and datasets. However, its functionality is primarily geared towards natural language processing.


% With the rapid proliferation of models and increasing complexity of HPC tasks, such an API can greatly streamline the deployment of new models, making it easier for researchers and practitioners to stay abreast of the latest developments in the field. In essence, the HPCFAIR API extends the \face API to the HPC domain, providing a more efficient and effective tool for managing, deploying, and evaluating HPC models.

\subsection{Design Overview}

%https://docs.google.com/drawings/d/18ZzDUncABpOFLPJXVhYJIgJlCdcmWsSh8yAKvwAp2LU/edit?usp=sharing
% \begin{figure}[h]
%     \centering
%     \includegraphics[width=12cm]{images/HPC-FAIR architecture.pdf}
%     \caption{Overview of the HPC-FAIR framework}
%     \label{fig:architecture}
% \end{figure}
\todo{HPCFAIR vs. HPC-FAIR in the figure ~\ref{fig:architecture}. datasets and evaluation harness (scripts) to generate leaderboards are also our new contributions, need to be highlighted.}
% \face, a platform embraced by the Machine Learning (ML) community, offers effortless access to a wide variety of pre-trained models and datasets. However, its primary focus is on Natural Language Processing (NLP), making it less suited for HPC tasks.

% To address the distinct needs of HPC tasks, we propose the HPC-FAIR API, built atop \face APIs. This approach allows us to harness the existing capabilities of the \face platform while incorporating additional features specifically designed for HPC tasks. Our design process begins by examining the typical workflows for HPC tasks.

% Pipelines offer intuitive interfaces for executing various HPC tasks using pre-trained models. This allows HPC researchers and developers, even those with limited ML experience, to rapidly harness the power of cutting-edge LLMs.

%Previous research ~\cite{verma2021hpcfair} introduces HPCFAIR, a modular and extensible framework to enable AI models. Here FAIR stands for Findable, Accessible, Interoperable, and Reproducible. In this work, we propose the HPC-FAIR framework following the FAIR principle.  HPC-FAIR is constructed on top of \face, a widely recognized machine learning platform. We broaden the applicability of HPC-FAIR, facilitating the integration of advanced language models into a wider array of HPC tasks.
Figure~\ref{fig:architecture} shows the overview of the HPC-FAIR framework. It is built on top of multiple internal machine learning components with \face-compatible APIs. Higher-level components provide concepts and interfaces to users, while middle or lower-level components provide implementation support. 
Example components introduced by HPC-FAIR include the inference pipelines for HPC-specific tasks, tokenization for programming languages, and relevant language models, datasets, and so on. %These additions enable HPC-FAIR to more effectively address the needs of AI researchers.

The term \textbf{workflow} describes a series of operations executed to accomplish a specific HPC ML goal. These goals may include using an existing model for a specific task, training a new model, evaluating and comparing multiple models, and so on. Each of these tasks is essential to the study of HPC with ML. Within the HPC-FAIR framework, we identify three primary workflows: inference pipeline, evaluation, and model training/fine-tuning.

Inference pipelines offer user-friendly interfaces for executing various HPC-related  tasks using pre-trained models. Example tasks including those related to HPC software analysis and optimizations, as well as question answering of HPC subjects. The pipelines allow HPC researchers and developers, even those with limited ML experience, to rapidly harness the power of cutting-edge LMs. 
%  We extended it to 

Evaluation workflows help one understand the strengths and limitations of the current LLM models in a standard, fair, and reproducible manner. Traditional metrics and processes may not be suitable or sufficient for HPC-specific tasks. To address this, new harnesses, compilers and tools, and metrics are needed. For example, specialized code processing tools are often needed to generate training raw datasets representing code information such as Abstract Syntax Tree, control-flow, and data-flow graphs.

Training/fine-tuning is an essential step for exploiting language models for HPC. Large Language Models (LLMs) are typically trained on a wide variety of internet text. However, only a select few of the most popular LLMs incorporate HPC programming languages within their training data. Even when a model has been trained with code, optimal performance for specialized tasks usually demands fine-tuning with user-specific data. To facilitate this, our HPC-FAIR framework provides APIs designed specifically for training and fine-tuning workflows.

% \subsection{HPC-FAIR Component}
% To achieve enable the HPC workflow, we need to expand the \face APIs. 

% 1. models and dataset. \face Hub is essentially a large repository for trained models that users can upload to and download from. The APIs for loading and using models are provided by the \face transformers library.  Table 4 lists the models on \face that have been specifically designed for code analysis and meet the requirements of HPC-FAIR workflows. 
\todo{The order of subsections can be arranged using a top-down fashion: framework overall, then test harness (script) benchmarking models (generating leaderboards), relevant tasks + pipelines, models, datasets, etc. from higher-level to low-level.  For each subsection, briefly mention relevant challenges and solutions, and examples.}
\subsection{Enabling Language Models}
% \face library provides APIs for users to upload/download models and datasets for NLP tasks from \face Hub, which is essentially a large repository. There are limited models in \face Hub that have been specifically designed for code analysis. 

Similar to NLP applications, HPC applications with language models encompass several essential components:
\subsubsection{Dataset}
\face Hub serves as a large repository for users to share NLP datasets. We process the HPC dataset \todo{what datasets? be specific.} and share it on \face. For example, users can load and inspect the pre-processed DataRaceBench dataset for HPC tasks evaluated in section \ref{sec-evaluation} using:
\begin{minted}
[
frame=lines,
framesep=2mm,
baselinestretch=1.2,
% bgcolor=LightGray,
fontsize=\footnotesize,
linenos
]{python}
from datasets import load_dataset
import pandas as pd
dataset = load_dataset("llnl/DRB-ML-Dataset")
# inspect the first example in the training set
print(dataset['train'][0]) 
df = pd.DataFrame(dataset['train'])
print(df.head())
\end{minted}

\subsubsection{Tokenizer}
% The data for HPC tasks are usually programming codes. The first step to process the dataset is data tokenization. A tokenizer is a component that splits the input code into smaller pieces called tokens. These tokens can be as short as one character or as long as one word, depending on the tokenizer. Most tokenizers on \face are designed for NLP tasks. 
\face tokenizers, primarily designed for Natural Language Processing (NLP) tasks, may encounter issues when directly applied to High-Performance Computing (HPC) tasks. For instance, consider the code in Listing \ref{ls:hello_w} processed with a BERT tokenizer, an NLP tokenizer. The generated tokens are \mintinline[breaklines, breakafter=]{bash}{['[cls]', '#', 'print', 'string', 'def', 'my', '', 'func', '():', '', 'print', '("', 'Hello', 'World', '")', '[sep]']}. Here, both the comment in line 1 and the function name  in line 2 are incorrectly split into multiple tokens, clearly illustrating the inadequacy of NLP tokenizers in parsing code correctly.

To overcome this limitation, we have developed the HPC-FAIR tokenizer, building it on top of the \face tokenizer. HPC-FAIR tokenizer is designed to handle all the pre-processing of code data required for a language model. It contains newly developed tokenizers like ast-tokenizer, a tokenizer parsing the code into tokens based on AST traversal. In line 4 of Listing \ref{ls:hpcfair-tokenizer}, the \textit{hpcfair-ast} tokenizer is loaded. The output of line 8 in Listing \ref{ls:hpcfair-tokenizer} is:  \mintinline[breaklines, breakafter=]{bash}{['[cls]', '# print string', 'def', 'my_func', '(', ')', ':', 'print', '(', '"Hello World"', ')', '[sep]']}. Comparing with the results of NLP tokenizers, the HPC-FAIR tokenizer parses the code more accurately.

\todo{This example~\ref{ls:hello_w} is too trivial, remove this Listing to save space}
 % We will conclude with more tokenizers designed for programming language in HPC-FAIR.
\begin{listing}[h]
\caption{A sample hello world python function. }
\label{ls:hello_w}
\begin{minted}
[
frame=lines,
framesep=2mm,
baselinestretch=1.2,
% bgcolor=LightGray,
fontsize=\footnotesize,
linenos
]{python}
# print string
def my_func():
    print("Hello World")
\end{minted}
\end{listing}

%\todo{What is the naming convention of python packages? from hpc-fair vs. from PC-FAIR:  \textit{Python package names should be all lowercase and avoid using underscores or hyphens.}}
\begin{listing}[h]
\caption{Example of using HPC-FAIR tokenizer.}
\label{ls:hpcfair-tokenizer}
\begin{minted}
[
frame=lines,
framesep=2mm,
baselinestretch=1.2,
% bgcolor=LightGray,
fontsize=\footnotesize,
linenos
]{python}
# import hpc-fair
from hpc-fair import AutoModel, AutoTokenizer
# load the tokenizer
tokenizer = AutoTokenizer.from_pretrained(hpcfair-ast)
# example code in Listing 1
code = "# print string\n def my_func(): \n print(\"Hello World\")"
tokens = tokenizer.tokenize(code)
print(tokens)
token_ids = tokenizer.encode(code, add_special_tokens=True)
# load the model from 1. \face Hub; 2. locally stored model checkpoint
model1 = AutoModel.from_pretrained("microsoft/codebert-base")
model2 = AutoModel.from_pretrained("/PATH/TO/CHECKPOINT")
inputs = torch.tensor([token_ids])
model = model1
outputs = model(inputs)
embeddings = outputs[0]
\end{minted}
\end{listing}






% \subsubsection{Embeddings}

\subsubsection{LM models}
Following the tokenization process, the data is prepared and suitable for input into the language models. As demonstrated in lines 11 and 12, there are two distinct methods of loading models.

The first method involves loading models from the \face Hub using a registered identifier. This approach provides easy access to a wide variety of pre-trained models and allows users to leverage the power of these models without having to train them from scratch.

The second method is to load models from a locally saved model checkpoint. This can be useful in cases where a user has fine-tuned a pre-existing model on their own specific data and saved the resulting model locally. This allows for greater flexibility and customization as users can continue from where they left off or reuse a particular model state.


% \subsubsection{Post-processing}

\subsection{Enabling pipelines}
% pick HPC tasks
% pick models
% fine-tune the model
% runt the inference pipelines
% evaluate the results
\begin{table}[htbp]
\centering
\caption{HPC Tasks for Programming Language Processing: Categories and Examples}
\label{tab-HPC-tasks}
\begin{tabular}{l|l|l}
\hline
%\multicolumn{3}{c}{\textbf{Category}} \\ \hline
\textbf{Code Analysis} & \textbf{Code Generation} & \textbf{Others} \\ \hline
Compiler Analysis & Code Completion & Test Case Generation \\
Algorithm Classification & Natural Language-to-Code & Code Search \\
Code Similarity Analysis & Code Translation & Question Answering \\
Documentation Generation & Code Repair & Decompilation \\
Parallelism Detection & Code Migration &  \\
Defect Detection & Compilation &  \\ \hline
\end{tabular}
\end{table}

Pipelines allows HPC researchers and developers, even those with limited ML experience, to rapidly harness the power of cutting-edge LLMs. Currently, \face provides pipelines for NLP and image processing tasks only. 
%\todo{HPCFAIR vs. HPC-FAIR: which one do you mean?  This paper focuses on HPC-FAIR. We only need to briefly mention previous HPCFAIR in related work. }
\subsubsection{HPC-FAIR tasks for LM application}
We first investigate HPC tasks suitable for using language models.
Table~\ref{tab-HPC-tasks} shows one way to categorize HPC-specific tasks related to programming language processing. The purpose here is not to provide a comprehensive taxonomy of all tasks but a starting point for common tasks we are interested in supporting in our framework. As the table shows, we roughly divide all tasks into three categories: code analysis, generation, and others. For each category, a few example tasks are given. Most tasks are self-explanatory by looking at their names. Each task may have further sub-tasks. For example, clone detection can be viewed as a specialized subtask under similarity analysis. 
\subsubsection{HPC-FAIR pipelines.}
Listing \ref{ls:hpcfair-tokenizer} demonstrates the steps involved in tokenizing code and processing it through models. This procedure can be complex for those with limited machine learning backgrounds. Challenges such as selecting the appropriate model, choosing the right tokenizer, and managing post-processing may demand substantial effort to address.

In an effort to streamline this process and provide user-friendly interfaces for executing a variety of HPC-related inference tasks using pre-trained models, we have extended the \face pipeline() API within HPC-FAIR. Currently, we have developed three pipelines: code similarity checking, parallelism detection, and OpenMP question answering. These pipelines are designed to significantly reduce the complexity of integrating machine learning models into HPC workflows.

Listing \ref{ls:pipepline} illustrates how to carry out inference tasks using HPC-FAIR. In line 3, the parallelism detection pipeline is loaded by passing the pipeline identifier to the \mintinline{python}{pipeline()} API. 
With just two lines of code, the pipeline API can generate results for parallelism detection tasks. Each pipeline task is equipped with predefined pre-trained models and tokenizers. This efficient design drastically reduces the complexity of these tasks, enabling users to obtain and evaluate results quickly and effortlessly.

\begin{listing}[h]
\caption{Example of using HPC-FAIR pipeline.}
\label{ls:pipepline}
\begin{minted}
[
frame=lines,
framesep=2mm,
baselinestretch=1.2,
% bgcolor=LightGray,
fontsize=\footnotesize,
linenos
]{python}
from hpc-fair import pipeline
# parallelism detection task
pipe = pipeline("parallelism-detection")
res = pipeline(code)
\end{minted}
\end{listing}
\todo{To save space, listings for API examples may be merged into a single listing, using line range start-end in the text to refer to different section. }


\todo[author=Leo, color=green]{It is recommended to list a few challenges (3-5) you encountered when designing/implementing the framework/API, and the corresponding solutions you used with some justifications.}

% \subsection{API Design}

\begin{table}[htbp]
	\centering
 \caption{API Comparison between HPC-FAIR and \face}
 \resizebox{\textwidth}{!}{
\begin{tabular}{|c|c|c|}
\hline
\textbf{Features} & \textbf{\face API} & \textbf{\ourtool API} \\ \hline
Audience & Generic  & HPC researchers and developers \\ \hline
Types of Models & NLP models (BERT, T5, etc.) & Text or Graph-based models \\ \hline
Sources of Models & Web-hosted, Local  & Web-hosted, Local, APIs  \\ \hline
Data Preprocessing & text, images  & Code  \\ \hline
Tokenization & NLP tokenizers (text-based) & Code in tokenized text, trees or graphs \\ \hline
Pipelines & For common NLP tasks & For common HPC PLP tasks \\ \hline
Training & Fine-tune training for NLP & Fine-tune training for PLP \\ \hline
Evaluation & Metrics and Harness for NLP & Metrics and Harness for HPC Tasks \\ \hline
Deployment & Web and mobile apps & large-scale HPC codebases with build systems \\ \hline
\end{tabular}
} % resizing
\label{tab-api-comparison}
\end{table}
\todo{\textcolor{yellow}{LangChain support loading models from HF. HF can work with langchain in  customized pipelines}Does HugginfFace support vector store and LangChain? If not, we should have a row in Table~\ref{tab-api-comparison} mentioning this: large input support: We support Vector Store, document split, code split, etc.}

\todo{need examples of our APIs: how they are used in python code}

% \subsection{Challenges}
% \begin{enumerate}
%     \item different model handling.
%     \item NLP tokenizer does not meet the needs for HPC PLP tasks.
%     \item limited dataset for QA.
%     \item extract interested texts only from the LLM outputs.
%     \item memory limitation: 1. better GPU 2. techniques like LLaMa-CPP
% \end{enumerate}


\todo{Where is the text about LangChain and benchmarking to generate leaderboard?}